% =====================================================================
%  TRƯỜNG ĐẠI HỌC CÔNG NGHỆ THÔNG TIN VÀ TRUYỀN THÔNG - ĐẠI HỌC THÁI NGUYÊN
%  BÁO CÁO TIỂU LUẬN HỌC PHẦN HỌC MÁY
%  Đề tài: OVERCOMING THE COLD-START PROBLEM IN DRUG-DRUG INTERACTION
%          PREDICTION VIA MULTI-SCENARIO ENSEMBLE LEARNING
%  Giảng viên hướng dẫn: TS. Nguyễn Tuấn Anh
%  Sinh viên thực hiện:
%    1. Trần Đặng Vương Quốc Long - MSV: DTC235200472 (Trưởng nhóm)
%    2. Bùi Minh Ảnh              - MSV: DTC235160004
%  Lớp: KHMT K22A - Khoa Công nghệ Thông tin - ICTU
%  Biên dịch hoàn hảo với: LaTeX thường (pdfLaTeX) hoặc XeLaTeX trên Overleaf
% =====================================================================
\documentclass[12pt,a4paper,oneside]{report}

% ---------------- Trang giấy ----------------
\usepackage[a4paper,left=3cm,right=2cm,top=2.5cm,bottom=2.5cm,
            headheight=16pt,headsep=0.8cm]{geometry}
\usepackage{type1cm}
\usepackage{iftex}

% ---------------- Hỗ trợ đa trình biên dịch (pdfLaTeX & XeLaTeX) ----------------
\ifPDFTeX
  % Dành cho LaTeX thường (pdfLaTeX - mặc định trên Overleaf)
  \usepackage[utf8]{vietnam}
  \usepackage{times}
\else
  % Dành cho XeLaTeX / LuaLaTeX
  \usepackage{fontspec}
  \setmainfont{Times New Roman}
  \setmonofont{Consolas}[Scale=0.9]
  \usepackage[vietnamese]{babel}
\fi

% Cỡ chữ thân bài 13pt (theo quy định chuẩn của bộ môn)
\makeatletter
\renewcommand\normalsize{%
  \@setfontsize\normalsize{13pt}{15.6pt}%
  \abovedisplayskip 12pt plus3pt minus7pt
  \belowdisplayskip \abovedisplayskip
  \abovedisplayshortskip 0pt plus3pt
  \belowdisplayshortskip 6.5pt plus3.5pt minus3pt}
\normalsize
\makeatother

% =====================================================================
%  THÔNG TIN BÁO CÁO
% =====================================================================
\newcommand{\truong}{TRƯỜNG ĐẠI HỌC CÔNG NGHỆ THÔNG TIN VÀ TRUYỀN THÔNG}
\newcommand{\khoa}{KHOA CÔNG NGHỆ THÔNG TIN}
\newcommand{\loaibaocao}{BÁO CÁO TIỂU LUẬN}
\newcommand{\hocphan}{Học máy}
\newcommand{\lophocphan}{KHMT K22A}
\newcommand{\hocky}{Học kỳ 1, năm học 2026--2027}
\newcommand{\tendetai}{OVERCOMING THE COLD-START PROBLEM IN DRUG-DRUG INTERACTION PREDICTION VIA MULTI-SCENARIO ENSEMBLE LEARNING}
\newcommand{\tenbaitoan}{dự đoán tương tác thuốc--thuốc}
\newcommand{\giangvien}{TS. Nguyễn Tuấn Anh}
\newcommand{\nhom}{Nhóm 01}
\newcommand{\svAten}{Trần Đặng Vương Quốc Long}
\newcommand{\svAma}{DTC235200472}
\newcommand{\svBten}{Bùi Minh Ảnh}
\newcommand{\svBma}{DTC235160004}
\newcommand{\diadiem}{Thái Nguyên}
\newcommand{\thoigian}{tháng 12 năm 2026}

\usepackage{setspace}
\setstretch{1.3}
\setlength{\parindent}{1cm}
\setlength{\parskip}{6pt}
\usepackage{indentfirst}

% ---------------- Gói hỗ trợ ----------------
\usepackage{amsmath,amssymb}
\usepackage{graphicx}
\graphicspath{{images/}{./}}
\usepackage{float}
\usepackage{array,booktabs,tabularx,multirow}
\renewcommand{\arraystretch}{1.25}
\usepackage[table]{xcolor}
\usepackage{tikz}
\usetikzlibrary{positioning,arrows.meta,calc,fit,backgrounds}
\usepackage{caption}
\captionsetup{labelsep=period,labelfont=bf,font=normalsize,skip=6pt}
\captionsetup[table]{position=top}
\usepackage{titlesec}
\usepackage{titletoc}
\usepackage[nottoc]{tocbibind}
\usepackage{enumitem}
\setlist{nosep,leftmargin=1.6cm}
\usepackage{listings}
\usepackage{needspace}
\usepackage{cite}
\usepackage{fancyhdr}
\usepackage[unicode]{hyperref}
\hypersetup{
  hidelinks,
  pdftitle={\tendetai},
  pdfauthor={\svAten, \svBten},
  pdfsubject={\loaibaocao\ học phần \hocphan},
  bookmarksnumbered=true,
  bookmarksopen=true,
  bookmarksopenlevel=0,
}
\newcommand{\bmchapname}{Chương}
\makeatletter
\newcommand*{\bm@prefix}{}
\ifdefined\Hy@writebookmark
  \let\bm@orig@writebookmark\Hy@writebookmark
  \renewcommand*{\Hy@writebookmark}[5]{%
    \ifnum#4=0 \def\bm@prefix{\bmchapname\space}\else\def\bm@prefix{}\fi
    \bm@orig@writebookmark{#1}{#2}{#3}{#4}{#5}}
\fi
\ifdefined\Hy@numberline
  \renewcommand*{\Hy@numberline}[1]{\bm@prefix#1. }
\fi
\makeatother

% ---------------- Màu & định dạng mã nguồn ----------------
\definecolor{codebg}{RGB}{247,247,247}
\definecolor{codekw}{RGB}{0,51,153}
\definecolor{codecm}{RGB}{0,128,0}
\definecolor{codestr}{RGB}{163,21,21}
\definecolor{guide}{RGB}{110,110,110}
\lstdefinestyle{python}{
  language=Python,
  basicstyle=\ttfamily\small\setstretch{1.0},
  backgroundcolor=\color{codebg},
  keywordstyle=\color{codekw}\bfseries,
  commentstyle=\color{codecm}\itshape,
  stringstyle=\color{codestr},
  numbers=left, numberstyle=\tiny\color{guide}, numbersep=8pt,
  frame=single, rulecolor=\color{black!25},
  breaklines=true, showstringspaces=false, tabsize=4,
  xleftmargin=1.2em, framexleftmargin=1.2em,
  aboveskip=8pt, belowskip=8pt, columns=fullflexible, keepspaces=true,
}
\lstset{style=python}
\renewcommand{\lstlistingname}{Đoạn mã}
\renewcommand{\lstlistlistingname}{DANH MỤC ĐOẠN MÃ}

% ---------------- Tiêu đề tiếng Việt chuẩn ----------------
\AtBeginDocument{%
  \renewcommand{\contentsname}{MỤC LỤC}%
  \renewcommand{\listfigurename}{DANH MỤC HÌNH ẢNH}%
  \renewcommand{\listtablename}{DANH MỤC BẢNG}%
  \renewcommand{\bibname}{TÀI LIỆU THAM KHẢO}%
  \renewcommand{\figurename}{Hình}%
  \renewcommand{\tablename}{Bảng}%
  \renewcommand{\chaptername}{Chương}%
  \renewcommand{\appendixname}{Phụ lục}%
}

% ---------------- Đầu trang & chân trang ----------------
\pagestyle{fancy}
\fancyhf{}
\fancyhead[L]{\small\itshape Báo cáo tiểu luận học phần \hocphan}
\fancyhead[R]{\small\itshape\nouppercase{\leftmark}}
\fancyfoot[C]{\thepage}
\renewcommand{\headrulewidth}{0.4pt}
\renewcommand{\chaptermark}[1]{\markboth{\chaptertitlename~\thechapter. #1}{}}
\fancypagestyle{plain}{\fancyhf{}\fancyfoot[C]{\thepage}\renewcommand{\headrulewidth}{0pt}}

% ---------------- Định dạng tiêu đề ----------------
\titleformat{\chapter}[display]
  {\filcenter\bfseries\fontsize{14}{18}\selectfont}
  {\MakeUppercase{\chaptertitlename}~\thechapter}{4pt}{\MakeUppercase}
\titlespacing*{\chapter}{0pt}{0pt}{18pt}
\titleformat{\section}
  {\bfseries\fontsize{14}{18}\selectfont}{\thesection.}{0.5em}{}
\titlespacing*{\section}{0pt}{12pt}{6pt}
\titleformat{\subsection}
  {\bfseries\fontsize{13}{17}\selectfont}{\thesubsection.}{0.5em}{}
\titlespacing*{\subsection}{0pt}{10pt}{6pt}
\setcounter{secnumdepth}{3}
\titleformat{\subsubsection}
  {\bfseries\itshape\fontsize{13}{17}\selectfont}{\thesubsubsection.}{0.5em}{}
\titlespacing*{\subsubsection}{0pt}{8pt}{4pt}

% ---------------- Mục lục & danh mục ----------------
\setcounter{tocdepth}{2}
\newcommand{\tocchapname}{Chương}
\titlecontents{chapter}[0pt]{\addvspace{6pt}\bfseries}
  {\tocchapname~\thecontentslabel.\enspace}{}{\titlerule*[0.6pc]{.}\contentspage}
\titlecontents{section}[1.5em]{}
  {\thecontentslabel.\enspace}{}{\titlerule*[0.6pc]{.}\contentspage}
\titlecontents{subsection}[3em]{}
  {\thecontentslabel.\enspace}{}{\titlerule*[0.6pc]{.}\contentspage}
\titlecontents{figure}[0pt]{}
  {\figurename~\thecontentslabel.\enspace}{}{\titlerule*[0.6pc]{.}\contentspage}
\titlecontents{table}[0pt]{}
  {\tablename~\thecontentslabel.\enspace}{}{\titlerule*[0.6pc]{.}\contentspage}

% ---------------- Lệnh tiện ích ----------------
\newcommand{\unnumberedchapter}[1]{%
  \chapter*{#1}\phantomsection\addcontentsline{toc}{chapter}{#1}\markboth{#1}{}}
\newcommand{\dongcham}[1]{\par\noindent\foreach \i in {1,...,#1}{\makebox[\linewidth]{\dotfill}\\[4pt]}}

\newcommand{\khunghinh}[3]{%
  \begin{tikzpicture}
    \draw[dashed,thick,black!40,fill=black!3] (0,0) rectangle (#1,#2);
    \node[text width=0.85\dimexpr#1\relax,align=center,font=\itshape\color{guide}]
      at ($ (0,0)!0.5!(#1,#2) $) {#3};
  \end{tikzpicture}}

% Lệnh chèn logo thông minh: Tự động tìm anh_nen.jpg, logo.jpg, logo.png
\newcommand{\chenlogo}{%
  \IfFileExists{anh_nen.jpg}{\includegraphics[width=3.2cm]{anh_nen.jpg}}{%
  \IfFileExists{anh_nen.png}{\includegraphics[width=3.2cm]{anh_nen.png}}{%
  \IfFileExists{logo.jpg}{\includegraphics[width=3.2cm]{logo.jpg}}{%
  \IfFileExists{logo.png}{\includegraphics[width=3.2cm]{logo.png}}{%
  \IfFileExists{images/logo.png}{\includegraphics[width=3.2cm]{images/logo.png}}{%
  \IfFileExists{images/anh_nen.jpg}{\includegraphics[width=3.2cm]{images/anh_nen.jpg}}{%
  \khunghinh{3.2cm}{3.2cm}{Logo\\trường}%
  }}}}}}%
}

% =====================================================================
\begin{document}

% =====================================================================
%  TRANG BÌA
% =====================================================================
\begin{titlepage}
  \begin{tikzpicture}[remember picture,overlay]
    \draw[line width=2.2pt]
      ($(current page.north west)+(2.4cm,-1.6cm)$) rectangle
      ($(current page.south east)+(-1.4cm,1.6cm)$);
    \draw[line width=0.6pt]
      ($(current page.north west)+(2.55cm,-1.75cm)$) rectangle
      ($(current page.south east)+(-1.55cm,1.75cm)$);
  \end{tikzpicture}\par\vspace*{-\baselineskip}
  \centering\setstretch{1.2}
  {\fontsize{13}{18}\selectfont\bfseries \truong\par}
  {\fontsize{14}{18}\selectfont\bfseries \khoa\par}
  \vspace{2pt}\rule{0.35\textwidth}{0.8pt}\par
  \vspace{0.8cm}
  \chenlogo\par
  \vspace{0.8cm}
  {\fontsize{20}{26}\selectfont\bfseries \loaibaocao\par}
  \vspace{4pt}
  {\fontsize{15}{19}\selectfont\bfseries HỌC PHẦN: \MakeUppercase{\hocphan}\par}
  \vspace{0.9cm}
  {\fontsize{13}{17}\selectfont ĐỀ TÀI:\par}
  \vspace{4pt}
  {\fontsize{15}{21}\selectfont\bfseries \tendetai\par}
  \vfill
  \begin{tabular}{@{}l@{\hspace{0.6cm}}l@{}}
    \textbf{Giảng viên hướng dẫn môn học:} & \giangvien \\
    \textbf{Lớp học phần:}                 & \lophocphan \\
    \textbf{Nhóm thực hiện:}               & \nhom \\
                                           & \svAten\ -- \svAma \\
                                           & \svBten\ -- \svBma \\
  \end{tabular}\par
  \vspace{1.2cm}
  {\bfseries \diadiem, \thoigian\par}
\end{titlepage}

% =====================================================================
%  PHẦN ĐẦU (đánh số trang La Mã)
% =====================================================================
\pagenumbering{roman}
\pagestyle{plain}

% ---------------- Nhận xét của giảng viên ----------------
\chapter*{NHẬN XÉT CỦA GIẢNG VIÊN}
\noindent\textbf{1. Về hình thức và bố cục báo cáo:}
\dongcham{3}
\noindent\textbf{2. Về nội dung và kết quả thực hiện:}
\dongcham{5}
\noindent\textbf{3. Về tinh thần, thái độ làm việc của nhóm:}
\dongcham{3}
\noindent\textbf{4. Điểm đánh giá:}
\begin{center}
  \begin{tabular}{|>{\centering\arraybackslash}p{1.2cm}|p{5.5cm}|>{\centering\arraybackslash}p{2.6cm}|>{\centering\arraybackslash}p{2.6cm}|}
    \hline
    \rowcolor{black!8}\textbf{STT} & \textbf{Họ và tên sinh viên} & \textbf{Điểm số} & \textbf{Điểm chữ} \\ \hline
    1 & \svAten & & \\ \hline
    2 & \svBten & & \\ \hline
  \end{tabular}
\end{center}
\vspace{0.6cm}
\hfill
\begin{minipage}{0.45\textwidth}
  \centering
  \textit{\diadiem, ngày \dots\ tháng \dots\ năm 2026}\\
  \textbf{Giảng viên}\\
  \textit{(Ký và ghi rõ họ tên)}\\[1.8cm]
  \giangvien
\end{minipage}

% ---------------- Lời cam đoan ----------------
\unnumberedchapter{LỜI CAM ĐOAN}

Chúng em xin cam đoan báo cáo tiểu luận ``\textit{\tendetai}'' là kết quả học tập và nghiên cứu thực nghiệm nghiêm túc, độc lập của nhóm dưới sự định hướng, hướng dẫn chuyên môn của giảng viên học phần \giangvien. Toàn bộ các kết quả thực nghiệm, số liệu thống kê khảo sát và bảng đối sánh mô hình trình bày trong báo cáo là trung thực, được kết xuất trực tiếp từ mã nguồn thực thi của nhóm và có thể tái lập hoàn toàn từ kho mã nguồn. Mọi tài liệu tham khảo, nền tảng lý thuyết và bộ dữ liệu kế thừa đều được trích dẫn nguồn gốc rõ ràng, đúng quy chuẩn học thuật.

\textbf{Khai báo sử dụng công cụ AI tạo sinh:} Nhóm có sử dụng công cụ AI hỗ trợ (ChatGPT / Claude / Antigravity) ở phạm vi: gợi ý dàn ý bố cục văn bản, rà soát chính tả, chuẩn hóa câu chữ và hỗ trợ định dạng trình bày bằng \LaTeX. Toàn bộ logic bài toán, thuật toán tiền xử lý dữ liệu, thiết lập 4 kịch bản kiểm soát rò rỉ dữ liệu, huấn luyện mô hình và phân tích chuyên môn đều do nhóm trực tiếp thực hiện, kiểm tra và làm chủ, sẵn sàng giải trình chi tiết trước giảng viên và hội đồng.

Nếu phát hiện có bất kỳ sự gian lận hoặc sao chép trái phép nào, chúng em xin hoàn toàn chịu mọi hình thức kỷ luật theo quy định của nhà trường.

\vspace{0.6cm}
\hfill
\begin{minipage}{0.65\textwidth}
  \centering
  \textit{\diadiem, ngày \dots\ tháng \dots\ năm 2026}\\
  \textbf{Nhóm sinh viên thực hiện}\\
  \textit{(Ký và ghi rõ họ tên)}\\[2cm]
  \begin{tabular}{@{}c@{\hspace{1.5cm}}c@{}}
    \svAten & \svBten \\
    \svAma  & \svBma  \\
  \end{tabular}
\end{minipage}

% ---------------- Lời cảm ơn ----------------
\unnumberedchapter{LỜI CẢM ƠN}

Để hoàn thành báo cáo tiểu luận môn học này, nhóm chúng em đã nhận được sự quan tâm, chỉ dẫn tận tình và giúp đỡ quý báu của quý thầy cô và bạn bè.

Trước hết, nhóm xin bày tỏ lòng biết ơn sâu sắc nhất tới thầy \textbf{\giangvien} -- giảng viên hướng dẫn học phần Học máy. Thầy đã tận tình truyền đạt những kiến thức nền tảng vững chắc về khoa học dữ liệu và học máy, đồng thời định hướng phương pháp luận khoa học, cách nhìn nhận bài toán khách quan, kiểm soát nghiêm ngặt hiện tượng rò rỉ dữ liệu và khích lệ tư duy nghiên cứu thực nghiệm của nhóm trong suốt học kỳ.

Nhóm xin chân thành cảm ơn Ban Giám hiệu cùng quý thầy cô Khoa Công nghệ Thông tin, Trường Đại học Công nghệ Thông tin và Truyền thông -- Đại học Thái Nguyên đã tạo điều kiện học tập, nghiên cứu và cung cấp môi trường học thuật thuận lợi nhất cho sinh viên. Nhóm cũng xin gửi lời cảm ơn tới các nhóm phát triển và duy trì các cơ sở dữ liệu y sinh DrugBank, SIDER và OFFSIDES đã công khai dữ liệu giá trị phục vụ mục đích nghiên cứu học thuật.

Mặc dù nhóm đã nỗ lực hết sức để hoàn thiện báo cáo một cách toàn diện và khoa học nhất, song do giới hạn về thời gian và trình độ, báo cáo chắc chắn không tránh khỏi những thiếu sót nhất định. Nhóm chúng em rất mong nhận được những góp ý, chỉ bảo quý báu của thầy cô để đề tài tiếp tục được hoàn thiện và nâng cao chất lượng.

Chúng em xin chân thành cảm ơn!

\vspace{0.4cm}
\hfill\textit{Nhóm sinh viên thực hiện: \svAten\ \& \svBten}

% ---------------- Bảng phân công ----------------
\unnumberedchapter{BẢNG PHÂN CÔNG NHIỆM VỤ}

\begin{center}
\begin{tabularx}{\textwidth}{|>{\centering\arraybackslash}p{1cm}|p{3.8cm}|X|>{\centering\arraybackslash}p{2.5cm}|}
  \hline
  \rowcolor{black!8}\textbf{STT} & \textbf{Họ và tên \& MSV} & \textbf{Nhiệm vụ phụ trách trong đề tài} & \textbf{Mức độ hoàn thành} \\ \hline
  1 & \svAten\newline \svAma\newline (Trưởng nhóm) &
    -- Khảo sát tài liệu nghiên cứu DDI và rà soát vấn đề rò rỉ dữ liệu ở các mô hình GNN/Ensemble\newline
    -- Thu thập, tích hợp dữ liệu 548 thuốc từ DrugBank 5.1.13 và xây dựng từ điển dữ liệu 10 kênh đặc trưng sinh y tế\newline
    -- Xây dựng thuật toán lọc Entropy và thuật toán hợp nhất mạng tương đồng (SNF)\newline
    -- Thiết lập 4 kịch bản phân chia dữ liệu chống rò rỉ (Warm-start, Strict Cold, Semi Cold, Mixed)\newline
    -- Soạn thảo báo cáo tiểu luận: Mở đầu, Chương 1, Chương 2 và tổng hợp toàn bộ báo cáo & 100\% \\ \hline
  2 & \svBten\newline \svBma &
    -- Thực hiện phân tích khám phá dữ liệu (EDA), phân tích phân bố nhãn, bậc tương tác và tương quan liên kênh\newline
    -- Xây dựng và tối ưu các mô hình phân loại: Logistic Regression, Random Forest, KNN, XGBoost và Soft Voting Ensemble\newline
    -- Triển khai mô hình cơ sở tham chiếu GCN Baseline và thực nghiệm 10 random seeds kiểm định độ vững chắc\newline
    -- Thực hiện thực nghiệm bóc tách thành phần (Ablation Study) và phân tích lỗi định tính\newline
    -- Soạn thảo báo cáo Chương 3, Kết luận và xây dựng Phụ lục A, B & 100\% \\ \hline
\end{tabularx}
\end{center}

% ---------------- Mục lục, danh mục ----------------
\begingroup\setstretch{1.1}\tableofcontents\endgroup
\listoffigures
\listoftables

% ---------------- Danh mục từ viết tắt ----------------
\unnumberedchapter{DANH MỤC TỪ VIẾT TẮT}

\begin{center}
\begin{tabularx}{\textwidth}{|>{\centering\arraybackslash}p{1.2cm}|p{2.6cm}|X|}
  \hline
  \rowcolor{black!8}\textbf{STT} & \textbf{Viết tắt} & \textbf{Ý nghĩa (Thuật ngữ tiếng Anh -- Tiếng Việt)} \\ \hline
  1 & ADE   & Adverse Drug Event -- Biến cố bất lợi của thuốc \\ \hline
  2 & ADR   & Adverse Drug Reaction -- Phản ứng có hại của thuốc \\ \hline
  3 & ATC   & Anatomical Therapeutic Chemical -- Hệ thống phân loại Giải phẫu -- Điều trị -- Hóa học \\ \hline
  4 & AUC   & Area Under the Receiver Operating Characteristic Curve \\ \hline
  5 & AUPR  & Area Under the Precision-Recall Curve -- Diện tích dưới đường cong PR \\ \hline
  6 & CV    & Cross-Validation -- Kiểm định chéo \\ \hline
  7 & DDI   & Drug--Drug Interaction -- Tương tác thuốc--thuốc \\ \hline
  8 & EDA   & Exploratory Data Analysis -- Phân tích khám phá dữ liệu \\ \hline
  9 & FN    & False Negative -- Âm tính giả (Bỏ sót tương tác thực tế) \\ \hline
  10 & FP    & False Positive -- Dương tính giả (Cảnh báo tương tác nhầm) \\ \hline
  11 & GCN   & Graph Convolutional Network -- Mạng tích chập đồ thị \\ \hline
  12 & GNN   & Graph Neural Network -- Mạng nơ-ron đồ thị \\ \hline
  13 & KNN   & K-Nearest Neighbors -- K láng giềng gần nhất \\ \hline
  14 & LR    & Logistic Regression -- Hồi quy Logistic \\ \hline
  15 & MeSH  & Medical Subject Headings -- Tiêu đề chủ đề y khoa \\ \hline
  16 & ML    & Machine Learning -- Học máy \\ \hline
  17 & RF    & Random Forest -- Rừng ngẫu nhiên \\ \hline
  18 & SMILES& Simplified Molecular Input Line Entry System -- Chuỗi biểu diễn cấu trúc hóa học \\ \hline
  19 & SNF   & Similarity Network Fusion -- Hợp nhất mạng tương đồng \\ \hline
  20 & TN    & True Negative -- Âm tính thật \\ \hline
  21 & TP    & True Positive -- Dương tính thật \\ \hline
  22 & XGBoost & eXtreme Gradient Boosting -- Tăng cường độ dốc cực đại \\ \hline
\end{tabularx}
\end{center}

% =====================================================================
%  PHẦN NỘI DUNG (đánh số trang Ả Rập)
% =====================================================================
\clearpage
\pagenumbering{arabic}
\pagestyle{fancy}

% =====================================================================
%  MỞ ĐẦU
% =====================================================================
\unnumberedchapter{MỞ ĐẦU}

\section*{1. Lý do chọn đề tài}
Trong y học điều trị hiện đại, tình trạng đa dược (polypharmacy) -- tức việc người bệnh phải sử dụng đồng thời từ hai loại thuốc trở lên để điều trị các bệnh mãn tính kết hợp (như tăng huyết áp, đái tháo đường, tim mạch) -- ngày càng trở nên phổ biến. Tuy nhiên, việc phối hợp nhiều hoạt chất tiềm ẩn rủi ro rất lớn phát sinh các tương tác thuốc--thuốc bất lợi (Drug--Drug Interactions -- DDI). DDI là nguyên nhân hàng đầu gây ra các phản ứng có hại của thuốc (Adverse Drug Reactions -- ADRs), làm suy giảm đáng kể hiệu quả điều trị hoặc thậm chí gây ngộ độc đe dọa trực tiếp đến tính mạng bệnh nhân \cite{pirmohamed2004adverse}.

Các thử nghiệm lâm sàng và nghiên cứu thực nghiệm trong phòng thí nghiệm (in vivo, in vitro) có độ tin cậy cao nhưng vô cùng tốn kém về tài chính, mất nhiều năm thực hiện và không thể kiểm thử hàng triệu cặp kết hợp thuốc khả dĩ trong thực tế. Do đó, việc ứng dụng học máy (Machine Learning) để dự đoán sớm các liên kết tương tác thuốc là một hướng tiếp cận cấp thiết và có giá trị thực tiễn rất cao.

Tuy nhiên, qua khảo sát các nghiên cứu hiện hành, nhóm nhận thấy một lỗ hổng phương pháp luận nghiêm trọng: đa số các công trình đánh giá mô hình bằng kỹ thuật chia ngẫu nhiên trên cấp độ cặp liên kết (random pair-level split). Cách chia này vô tình tạo ra hiện tượng \textbf{rò rỉ dữ liệu cấu trúc (topological data leakage)}: các thuốc trong tập kiểm tra đã từng xuất hiện ở tập huấn luyện. Dẫn đến việc mô hình chỉ học thuộc lòng cấu trúc đồ thị quen thuộc thay vì học được quy luật sinh hóa thực sự. Khi đối mặt với các loại thuốc mới chưa từng có lịch sử kê đơn trong mạng lưới (bài toán khởi động lạnh -- Cold-Start), hiệu năng của các mô hình này bị sụp đổ nghiêm trọng. 

Nhận thức rõ thách thức đó, nhóm đã lựa chọn đề tài: \textbf{``\tendetai''}.

\section*{2. Mục tiêu và nhiệm vụ của đề tài}
\textbf{Mục tiêu tổng quát:} Thiết lập một quy trình học máy minh bạch, khoa học từ dữ liệu đến mô hình; đề xuất khung đánh giá đa kịch bản (Multi-Scenario Framework) kiểm soát triệt để hiện tượng rò rỉ dữ liệu; đồng thời xây dựng mô hình Soft Voting Ensemble kết hợp đa nguồn đặc trưng sinh học nhằm dự đoán chính xác và ổn định tương tác thuốc trong các bối cảnh khởi động lạnh khắt khe.

\textbf{Nhiệm vụ cụ thể:}
\begin{enumerate}[label=\alph*)]
  \item Chuẩn hóa phát biểu bài toán dự đoán liên kết DDI dưới dạng bài toán phân loại nhị phân có giám sát; lượng hóa các mục tiêu thực tế thành hệ chỉ số kỹ thuật khách quan.
  \item Tích hợp dữ liệu 548 thuốc chuẩn từ DrugBank 5.1.13, SIDER và OFFSIDES; xây dựng không gian biểu diễn đa góc nhìn gồm 10 kênh đặc trưng sinh y tế.
  \item Áp dụng kỹ thuật lọc đặc trưng dựa trên Entropy và thuật toán hợp nhất mạng tương đồng (Similarity Network Fusion -- SNF) để tạo ra biểu diễn tổng hợp chất lượng cao.
  \item Thiết kế và triển khai 4 kịch bản phân chia dữ liệu chống rò rỉ: Warm-start, Strict Cold-start (New $\times$ New), Semi Cold-start (Old $\times$ New) và Mixed.
  \item Huấn luyện và so sánh 5 thuật toán học máy phân loại đại diện (Logistic Regression, Random Forest, KNN, XGBoost, Soft Voting Ensemble) và mô hình học sâu đồ thị GCN Baseline; thực hiện kiểm định độ vững chắc qua 10 random seeds; bóc tách đóng góp thành phần (Ablation Study) và phân tích lỗi định tính.
\end{enumerate}

\section*{3. Đối tượng và phạm vi nghiên cứu}
\begin{itemize}
  \item \textbf{Đối tượng nghiên cứu:} Tập dữ liệu 548 thuốc phổ biến, đồ thị 48.584 tương tác DDI đã được xác thực lâm sàng (tương ứng 97.168 phần tử khác không trong ma trận kề đối xứng); 10 nguồn thông tin sinh học (SMILES, Protein đích, Enzyme, Transporter, Pathway, Chỉ định điều trị, Tác dụng phụ, Biến cố bất lợi, Mã ATC, Mã MeSH) và các thuật toán học máy phân loại.
  \item \textbf{Phạm vi nghiên cứu:} Bài toán phân loại nhị phân dự đoán khả năng tương tác của cặp thuốc trên 4 kịch bản phân chia dữ liệu kiểm soát rò rỉ; nghiên cứu chưa mở rộng sang dự đoán loại tương tác đa nhãn (multi-typed DDI).
\end{itemize}

\section*{4. Phương pháp nghiên cứu}
Đề tài kết hợp phương pháp nghiên cứu lý thuyết và thực nghiệm trên nền tảng Python (Scikit-learn, XGBoost, RDKit, PyTorch Geometric, cuML). Dữ liệu được tiền xử lý chuẩn hóa, áp dụng chiến lược phân tách tập thuốc rời rạc (drug-disjoint node split) nhằm ngăn chặn tuyệt đối việc rò rỉ thực thể thuốc giữa tập huấn luyện và kiểm tra.

\section*{5. Bố cục của báo cáo tiểu luận}
Báo cáo được kết cấu thành 3 chương chính:
\begin{itemize}
  \item \textbf{Chương 1. Tổng quan về học máy và bài toán \tenbaitoan:} Trình bày cơ sở lý thuyết học máy, bài toán dự đoán liên kết đồ thị, tổng quan các nghiên cứu liên quan, biểu diễn đặc trưng, thuật toán SNF và các độ đo đánh giá.
  \item \textbf{Chương 2. Bài toán và phương pháp xử lý dữ liệu:} Chi tiết hóa bài toán, quy trình xử lý dữ liệu tổng quát, xây dựng 10 kênh đặc trưng sinh học, tinh lọc Entropy, hợp nhất SNF, thiết lập 4 kịch bản chia tập chống rò rỉ và phân tích khám phá dữ liệu (EDA).
  \item \textbf{Chương 3. Thực nghiệm và đánh giá:} Thiết lập thí nghiệm, mô tả các mô hình, phân tích kết quả đối sánh 5 mô hình trên 4 kịch bản, đánh giá độ vững chắc so với GCN Baseline qua 10-seed, phân tích lỗi định tính và thực nghiệm bóc tách thành phần (Ablation Study).
\end{itemize}

% =====================================================================
%  CHƯƠNG 1
% =====================================================================
\chapter{Tổng quan về học máy và bài toán \tenbaitoan}
\label{chap:tongquan}

\section{Tổng quan về học máy}

\subsection{Khái niệm học máy}
Học máy (Machine Learning -- ML) là một phân ngành cốt lõi của trí tuệ nhân tạo, tập trung nghiên cứu và phát triển các thuật toán cho phép hệ thống máy tính tự động đúc rút quy luật từ dữ liệu thực nghiệm thay vì phụ thuộc hoàn toàn vào các quy tắc lập trình thủ công \cite{mitchell1997}. Theo định nghĩa kinh điển của Tom Mitchell, một chương trình máy tính được gọi là học từ kinh nghiệm $E$ đối với một tập nhiệm vụ $T$ và độ đo hiệu năng $P$, nếu hiệu năng thực hiện nhiệm vụ thuộc $T$ đo bằng $P$ được cải thiện theo kinh nghiệm $E$.

Trong khuôn khổ đề tài này: nhiệm vụ $T$ là dự đoán xác suất tương tác giữa một cặp thuốc bất kỳ; kinh nghiệm $E$ là tập các cặp thuốc đã biết có hoặc chưa ghi nhận tương tác cùng các đặc trưng sinh học trong DrugBank; độ đo hiệu năng $P$ là các chỉ số khách quan như ROC-AUC, AUPR và F1-score.

\subsection{Phân loại các phương pháp học máy}
Dựa trên hình thức dữ liệu và tín hiệu phản hồi, học máy được chia thành 4 nhóm chính:
\begin{itemize}
  \item \textbf{Học có giám sát (Supervised Learning):} Dữ liệu huấn luyện gồm các cặp mẫu và nhãn $(x_i, y_i)$. Mục tiêu là học hàm ánh xạ $f: x \mapsto y$.
  \item \textbf{Học không giám sát (Unsupervised Learning):} Dữ liệu không có nhãn; mục tiêu là phát hiện các cấu trúc tiềm ẩn, phân cụm hoặc giảm chiều dữ liệu (như thuật toán SNF, phân cụm đồ thị).
  \item \textbf{Học bán giám sát (Semi-supervised Learning):} Kết hợp lượng nhỏ dữ liệu có nhãn với lượng lớn dữ liệu chưa gán nhãn.
  \item \textbf{Học tăng cường (Reinforcement Learning):} Tác tử học cách ra quyết định thông qua tương tác thử--sai với môi trường để tối đa hóa phần thưởng.
\end{itemize}

\subsection{Quy trình xây dựng dự án học máy và vấn đề quá khớp}
Một dự án học máy tiêu chuẩn trải qua chu trình lặp: Phát biểu bài toán $\rightarrow$ Thu thập dữ liệu $\rightarrow$ Tiền xử lý \& Khảo sát EDA $\rightarrow$ Huấn luyện mô hình $\rightarrow$ Đánh giá $\rightarrow$ Tối ưu hóa \& Triển khai. Hai rủi ro lớn nhất cần kiểm soát là:
\begin{itemize}
  \item \textbf{Chưa khớp (Underfitting):} Mô hình quá đơn giản, không học được quy luật của dữ liệu (độ chệch -- bias cao).
  \item \textbf{Quá khớp (Overfitting):} Mô hình ghi nhớ cả nhiễu của tập huấn luyện, đạt độ chính xác cực cao trên tập train nhưng thất bại khi dự đoán dữ liệu mới (phương sai -- variance cao).
\end{itemize}

\section{Bài toán phân loại có giám sát và thách thức mất cân bằng lớp}
Bài toán phân loại nhị phân tìm kiếm siêu phẳng hoặc hàm quyết định $f(\mathbf{x}): \mathbb{R}^d \rightarrow [0, 1]$ cực tiểu hóa rủi ro thực nghiệm có điều chuẩn:
\begin{equation}
  \hat{f} = \arg\min_{f \in \mathcal{F}} \frac{1}{n}\sum_{i=1}^{n} \mathcal{L}(y_i, f(\mathbf{x}_i)) + \lambda \Omega(f).
\end{equation}
Trong bài toán tương tác thuốc, số lượng cặp không tương tác (hoặc chưa ghi nhận) áp đảo hoàn toàn so với số cặp tương tác đã biết. Hiện tượng mất cân bằng lớp (class imbalance) khiến độ chính xác Accuracy trở nên sai lệch: một mô hình tầm thường dự đoán mọi cặp đều là 0 vẫn có thể đạt Accuracy trên 80\% nhưng hoàn toàn vô giá trị về mặt y học. Do đó, việc áp dụng kỹ thuật lấy mẫu cân bằng (Balanced Negative Sampling 1:1) và ưu tiên chỉ số AUPR là bắt buộc.

\section{Học máy trên dữ liệu đồ thị và bài toán dự đoán liên kết}
Đồ thị vô hướng $G = (V, E)$ biểu diễn mạng lưới tương tác thuốc, trong đó tập đỉnh $V$ là các loại thuốc ($|V| = N = 548$) và tập cạnh $E$ là các tương tác thực tế ($|E| = M = 48.584$). Bài toán dự đoán liên kết (Link Prediction) ước lượng khả năng tồn tại cạnh giữa hai đỉnh $u$ và $v$. Tiếp cận theo hướng học máy có giám sát chuyển bài toán này thành việc xây dựng vector đặc trưng $\mathbf{x}_{uv}$ cho mỗi cặp đỉnh và dự đoán nhãn $y_{uv} \in \{0, 1\}$.

\section{Bài toán dự đoán tương tác thuốc--thuốc (DDI) và các nghiên cứu liên quan}
\begin{itemize}
  \item \textbf{Tiếp cận dựa trên cấu trúc phân tử:} Vilar và cộng sự \cite{vilar2014} sử dụng độ tương đồng cấu trúc 2D. Hạn chế lớn là các thuốc có cấu trúc hóa học khác xa nhau nhưng có thể cùng tương tác qua cùng một thụ thể hoặc enzyme chuyển hóa.
  \item \textbf{Tiếp cận tích hợp đa nguồn sinh học:} Rohani và Eslahchi \cite{rohani2019} đề xuất mô hình NDD, tích hợp các kênh enzyme, đích sinh học, con đường chuyển hóa và chỉ định bệnh học bằng mạng nơ-ron sâu. Tuy nhiên, mô hình sử dụng kiểm định chéo ngẫu nhiên trên các cặp thuốc, gây rò rỉ dữ liệu nghiêm trọng.
  \item \textbf{Tiếp cận bằng mạng nơ-ron đồ thị (GNN):} Zitnik và cộng sự \cite{zitnik2018} phát triển Decagon trên đồ thị đa quan hệ thuốc--protein. Tuy nhiên, các nghiên cứu gần đây (2024--2026) \cite{langground2026, fgddi2026, mrlfddi2026} chỉ ra rằng GNN phụ thuộc nặng nề vào cài đặt transductive và bị sụp đổ hiệu năng nghiêm trọng khi gặp thuốc hoàn toàn mới trong kịch bản inductive/cold-start.
\end{itemize}

\section{Các kỹ thuật biểu diễn đặc trưng cho cặp thuốc}
Để mô tả toàn diện bản chất sinh hóa học của một cặp thuốc, đề tài xây dựng 10 không gian đặc trưng bổ trợ:
\begin{enumerate}
  \item \textbf{Cấu trúc hóa học (Chem):} Chuỗi SMILES được chuyển thành vân tay phân tử Morgan (Morgan Fingerprints, bán kính $r=2$, độ dài 1024 bits) bằng thư viện RDKit. Độ tương đồng giữa hai tập bit $A$ và $B$ được đo bằng hệ số Tanimoto:
  \begin{equation}
    \text{Tanimoto}(A, B) = \frac{|A \cap B|}{|A \cup B|}.
  \end{equation}
  \item \textbf{Chỉ định điều trị (Indication):} Văn bản y khoa mô tả bệnh học được mã hóa trọng số TF-IDF và tính độ tương đồng bằng Cosine similarity:
  \begin{equation}
    \text{Cosine}(\mathbf{u}, \mathbf{v}) = \frac{\mathbf{u} \cdot \mathbf{v}}{\|\mathbf{u}\| \|\mathbf{v}\|}.
  \end{equation}
  \item \textbf{8 kênh liên kết sinh học còn lại (Target, Enzyme, Pathway, Transporter, Side Effect, Offside Effect, ATC, MeSH):} Được biểu diễn dưới dạng vector nhị phân (one-hot binary profile) và đo độ tương đồng bằng chỉ số Jaccard:
  \begin{equation}
    J(A, B) = \frac{|A \cap B|}{|A \cup B|}.
  \end{equation}
  Khi thuốc thiếu thuộc tính, gán vector 0 và giá trị độ tương đồng bằng 0 để bảo toàn tính khách quan.
\end{enumerate}

\section{Phương pháp lọc đặc trưng bằng Entropy và Hợp nhất mạng (SNF)}
\begin{itemize}
  \item \textbf{Lọc đặc trưng bằng Entropy:} Để loại bỏ các đặc trưng chứa nhiều nhiễu, Entropy của từng vector đặc trưng được tính: $E = -\sum p \log p$. Những đặc trưng có tỷ lệ entropy $\gamma = E / E_{\max} < 0{,}85$ sẽ bị loại bỏ.
  \item \textbf{Hợp nhất mạng tương đồng (SNF):} Áp dụng thuật toán SNF \cite{wang2014} với các tham số chuẩn $K=20$ (số láng giềng) và $t=20$ (số bước lặp). SNF khuếch tán thông tin giữa 10 ma trận tương đồng thông qua phép nhân ma trận lặp:
  \begin{equation}
    P^{(t+1)}_i = S_i \times \left(\frac{1}{m-1} \sum_{j \neq i} P^{(t)}_j\right) \times S_i^{\top},
  \end{equation}
  tạo nên một ma trận hợp nhất thống nhất, giàu thông tin và giảm thiểu khoảng trống dữ liệu thiếu.
\end{itemize}

\section{Các độ đo đánh giá mô hình}
Để đánh giá toàn diện bài toán phân loại nhị phân trên tập dữ liệu mất cân bằng, nhóm sử dụng ma trận nhầm lẫn gồm $TP, FP, TN, FN$ và các độ đo:
\begin{equation}
  \text{Precision} = \frac{TP}{TP + FP}, \quad \text{Recall} = \frac{TP}{TP + FN}, \quad F_1 = \frac{2 \cdot \text{Precision} \cdot \text{Recall}}{\text{Precision} + \text{Recall}}.
\end{equation}
Đặc biệt, hai độ đo không phụ thuộc ngưỡng quyết định:
\begin{itemize}
  \item \textbf{ROC-AUC:} Diện tích dưới đường cong ROC đo lường năng lực phân biệt tổng quát giữa hai lớp.
  \item \textbf{PR-AUC (AUPR):} Diện tích dưới đường cong Precision-Recall. Trong bài toán y sinh, AUPR là độ đo cốt lõi quan trọng nhất vì nó phản ánh trực tiếp năng lực phát hiện các tương tác thuốc có hại trong bối cảnh lớp dương chiếm tỷ lệ nhỏ.
\end{itemize}

% =====================================================================
%  CHƯƠNG 2
% =====================================================================
\chapter{Bài toán và phương pháp xử lý dữ liệu}
\label{chap:phuongphap}

\section{Phát biểu bài toán và những thách thức cốt lõi}
\subsection{Phát biểu bài toán}
Cho tập hợp 548 thuốc $D = \{d_1, d_2, \dots, d_{548}\}$ và đồ thị tương tác DDI đã biết $G = (V, E)$. Với mỗi cặp thuốc $(d_i, d_j)$ với $i < j$, bài toán yêu cầu xây dựng vector biểu diễn $\mathbf{x}_{ij}$ từ 10 kênh sinh học và dự đoán xác suất tương tác $\hat{y}_{ij} = P(y=1 \mid \mathbf{x}_{ij}) \in [0, 1]$.

\subsection{Các thách thức kỹ thuật}
\begin{enumerate}
  \item \textbf{Rò rỉ dữ liệu cấu trúc (Topological Leakage):} Khi chia tập theo cặp ngẫu nhiên, các đỉnh test xuất hiện trong tập train khiến mô hình chỉ ghi nhớ cấu trúc đồ thị.
  \item \textbf{Bài toán khởi động lạnh (Cold-Start):} Khi một thuốc hoàn toàn mới được cấp phép lâm sàng, thuốc đó chưa hề có bất kỳ liên kết DDI nào. Mô hình buộc phải suy diễn dựa hoàn toàn trên đặc trưng sinh hóa học.
  \item \textbf{Bản chất nhãn PU (Positive-Unlabeled):} Các cặp nhãn 0 thực chất là ``chưa quan sát'' (unobserved), có thể chứa các tương tác lâm sàng chưa được phát hiện.
\end{enumerate}

\section{Dữ liệu thực nghiệm và quy mô sau xử lý}
Dữ liệu được chuẩn hóa từ danh sách 548 thuốc của Rohani và cộng sự \cite{rohani2019}, cập nhật thông tin từ DrugBank phiên bản 5.1.13, SIDER và OFFSIDES.
Tổng số cặp vô hướng khả dĩ trong không gian 548 thuốc là:
\begin{equation}
  N_{\text{pairs}} = \frac{548 \times 547}{2} = 149.878\text{ cặp}.
\end{equation}
Trong đó, số cặp tương tác dương tính ($y=1$) là 48.584 cặp (tương ứng 97.168 phần tử khác không trong ma trận đối xứng). Số cặp chưa ghi nhận tương tác ($y=0$) là 101.294 cặp. Để kiểm soát mất cân bằng lớp và tối ưu hóa hàm mất mát, nhóm áp dụng kỹ thuật Balanced Negative Sampling lấy đúng 48.584 mẫu âm, tạo thành tập dữ liệu cân bằng 97.168 mẫu với tỷ lệ 1:1.

\section{Quy trình xử lý dữ liệu tổng quát}
Quy trình tiền xử lý và xây dựng đặc trưng được chuẩn hóa theo 4 giai đoạn mạch lạc:
\begin{enumerate}
  \item \textbf{Trích xuất \& Tiền xử lý đa nguồn:} Thu thập 10 kênh đặc trưng, xử lý dữ liệu thiếu bằng cách gán vector 0 và tính 10 ma trận tương đồng đối xứng kích thước $548 \times 548$.
  \item \textbf{Lọc Entropy \& Hợp nhất SNF:} Lọc bỏ các kênh đặc trưng có tỷ lệ entropy $\gamma < 0{,}85$, sau đó áp dụng thuật toán SNF lặp 20 bước với $K=20$ để tạo ma trận tương đồng hợp nhất.
  \item \textbf{Phân chia tập thuốc rời rạc (Drug-Disjoint Splitting):} Tách danh sách 548 thuốc thành tập huấn luyện $D_{\text{train}}$ và tập kiểm tra $D_{\text{test}}$ với điều kiện nghiêm ngặt $D_{\text{train}} \cap D_{\text{test}} = \emptyset$.
  \item \textbf{Tạo lập 4 kịch bản đánh giá:} Sinh các cặp tương tác dương và lấy mẫu âm cân bằng 1:1 theo đúng cấu hình của từng kịch bản.
\end{enumerate}

\begin{figure}[H]
  \centering
  \resizebox{\textwidth}{!}{%
  \begin{tikzpicture}[
    font=\small,
    node distance=0.55cm,
    block/.style={draw=black!70, line width=0.8pt, rounded corners=4pt, fill=black!4, align=center, minimum height=1.3cm, text width=2.8cm},
    data/.style={draw=black!70, line width=0.8pt, rounded corners=4pt, fill=black!10, align=center, minimum height=1.1cm, text width=2.4cm},
    arrow/.style={-{Stealth[length=2.5mm]}, thick, draw=black!70}
  ]
    \node[block] (f) {\textbf{Đặc trưng đa nguồn}\\\scriptsize 10 kênh sinh y tế\\(SMILES, Target, ADE...)};
    \node[block, right=of f] (e) {\textbf{Lọc Entropy}\\\scriptsize Tinh lọc giảm nhiễu\\($\gamma \ge 0{,}85$)};
    \node[block, right=of e] (s) {\textbf{Hợp nhất SNF}\\\scriptsize Hợp nhất mạng\\($K=20, t=20$)};
    \node[block, right=of s] (c) {\textbf{Soft Voting}\\\scriptsize RF, KNN, XGB, LR\\(Trung bình xác suất)};
    \node[data, right=of c] (o) {\textbf{Dự đoán DDI}\\\scriptsize $\hat{y} \in \{0, 1\}$\\(Ngưỡng $0{,}5$)};

    \draw[arrow] (f) -- (e);
    \draw[arrow] (e) -- (s);
    \draw[arrow] (s) -- (c);
    \draw[arrow] (c) -- (o);
  \end{tikzpicture}%
  }
  \caption{Sơ đồ quy trình dự đoán tương tác thuốc--thuốc tích hợp đa nguồn đặc trưng và mô hình Soft Voting Ensemble}
  \label{fig:pipeline_overview}
\end{figure}

\section{Thiết lập 4 kịch bản đánh giá chống rò rỉ dữ liệu}
Đây là đóng góp phương pháp luận quan trọng nhất của đề tài, mô phỏng từ mức độ dễ đến các tình huống thử thách khắt khe nhất trong y tế thực tế:
\begin{itemize}
  \item \textbf{Case 1: Warm-start (Random pair split):} Cả hai thuốc trong cặp kiểm tra đều đã từng xuất hiện trong tập huấn luyện ($u, v \in D_{\text{train}}$). Mô hình khai thác tối đa cấu trúc đồ thị quen thuộc.
  \item \textbf{Case 2: Strict Cold-start (New $\times$ New):} Cả hai thuốc trong cặp kiểm tra đều hoàn toàn mới và chưa từng xuất hiện ở tập huấn luyện ($u, v \in D_{\text{test}}$). Đánh giá năng lực suy luận sinh học thuần túy không rò rỉ cấu trúc.
  \item \textbf{Case 3: Semi Cold-start (Old $\times$ New):} Một thuốc đã biết trong tập huấn luyện kết hợp với một thuốc hoàn toàn mới ($u \in D_{\text{train}}, v \in D_{\text{test}}$), phản ánh đúng bối cảnh kê đơn phối hợp thuốc mới với các thuốc lưu hành lâu năm.
  \item \textbf{Case 4: Mixed Scenario:} Kịch bản hỗn hợp gộp cả Case 2 và Case 3, phản ánh môi trường lâm sàng đa dạng tại các bệnh viện.
\end{itemize}

\begin{figure}[H]
  \centering
  \resizebox{\textwidth}{!}{%
  \begin{tikzpicture}[
    font=\small,
    box/.style={draw=black!60, line width=0.8pt, rounded corners=4pt, fill=black!3, minimum width=3.0cm, minimum height=3.5cm, align=center},
    header/.style={font=\bfseries\small, align=center},
    seen/.style={circle, draw=black!80, fill=black!25, minimum size=0.65cm, inner sep=0pt, font=\scriptsize\bfseries},
    unseen/.style={circle, draw=black!80, fill=black!5, minimum size=0.65cm, inner sep=0pt, font=\scriptsize\bfseries},
    edge/.style={thick, draw=black!70}
  ]
    % Case 1
    \node[box] (b1) at (0,0) {};
    \node[header, below=4pt of b1.north] {Case 1: Warm};
    \node[seen] (u1) at (-0.6, 0.3) {$u$};
    \node[seen] (v1) at (0.6, 0.3) {$v$};
    \draw[edge] (u1) -- (v1);
    \node[below=0.2cm of b1.center, font=\scriptsize, align=center] {$u, v \in D_{\text{train}}$\\\textit{(Đã biết cả hai)}};

    % Case 2
    \node[box] (b2) at (3.6,0) {};
    \node[header, below=4pt of b2.north] {Case 2: Strict};
    \node[unseen] (u2) at (3.0, 0.3) {$u$};
    \node[unseen] (v2) at (4.2, 0.3) {$v$};
    \draw[edge, dashed] (u2) -- (v2);
    \node[below=0.2cm of b2.center, font=\scriptsize, align=center] {$u, v \in D_{\text{test}}$\\\textit{(Cả hai đều mới)}};

    % Case 3
    \node[box] (b3) at (7.2,0) {};
    \node[header, below=4pt of b3.north] {Case 3: Semi};
    \node[seen] (u3) at (6.6, 0.3) {$u$};
    \node[unseen] (v3) at (7.8, 0.3) {$v$};
    \draw[edge] (u3) -- (v3);
    \node[below=0.2cm of b3.center, font=\scriptsize, align=center] {$u \in D_{\text{train}}, v \in D_{\text{test}}$\\\textit{(Một cũ -- một mới)}};

    % Case 4
    \node[box] (b4) at (10.8,0) {};
    \node[header, below=4pt of b4.north] {Case 4: Mixed};
    \node[above=0.2cm of b4.center, font=\scriptsize, align=center] {\textbf{Hợp nhất test set}\\\textit{Case 2 + Case 3}};
    \node[below=0.2cm of b4.center, font=\scriptsize, align=center] {Mô phỏng luồng\\lâm sàng hỗn hợp};

    % Legend
    \node[draw=black!40, rounded corners=3pt, fill=white, below=0.35cm of $(b2)!0.5!(b3)$, font=\scriptsize, inner sep=4pt] {
      \textbf{Chú giải:}\quad
      \begin{tikzpicture}[baseline=-0.5ex]
        \node[seen, minimum size=0.42cm] {};
      \end{tikzpicture}
      Thuốc trong tập Train ($D_{\text{train}}$)\qquad
      \begin{tikzpicture}[baseline=-0.5ex]
        \node[unseen, minimum size=0.42cm] {};
      \end{tikzpicture}
      Thuốc mới trong tập Test ($D_{\text{test}}$)
    };
  \end{tikzpicture}%
  }
  \caption{Minh họa 4 kịch bản phân chia dữ liệu đánh giá độ vững chắc và kiểm soát rò rỉ}
  \label{fig:scenarios_overview}
\end{figure}

Bảng~\ref{tab:scenarios} trình bày chi tiết số lượng mẫu và mật độ đồ thị của từng kịch bản theo công bố thực nghiệm của bài báo.

\begin{table}[H]
  \centering
  \caption{Đặc tính phân chia dữ liệu của 4 kịch bản thực nghiệm chống rò rỉ}
  \label{tab:scenarios}
  \begin{tabularx}{\textwidth}{l c c >{\centering\arraybackslash}p{3.2cm} >{\centering\arraybackslash}p{2.6cm}}
    \toprule
    \textbf{Kịch bản đánh giá} & \textbf{Thuốc $u$} & \textbf{Thuốc $v$} & \textbf{Số cặp Dương / Âm} & \textbf{Mật độ đồ thị (\%)} \\
    \midrule
    Case 1 (Warm-start)       & $u \in D_{\text{train}}$ & $v \in D_{\text{train}}$ & 9.716 / 9.716   & 32,42 \\
    Case 2 (Strict Cold-start) & $u \in D_{\text{test}}$  & $v \in D_{\text{test}}$  & 1.728 / 1.728   & 28,82 \\
    Case 3 (Semi Cold-start)   & $u \in D_{\text{train}}$ & $v \in D_{\text{test}}$  & 15.077 / 15.077 & 31,29 \\
    Case 4 (Mixed Scenario)   & \multicolumn{2}{c}{Hợp nhất của Case 2 \& Case 3}  & 16.805 / 16.805 & 31,02 \\
    \bottomrule
  \end{tabularx}
\end{table}

\section{Khảo sát và phân tích khám phá dữ liệu (EDA)}
\subsection{Từ điển dữ liệu 10 kênh đặc trưng}
Từ điển dữ liệu tóm tắt 10 không gian đặc trưng sinh y học của 548 thuốc được mô tả trong Bảng~\ref{tab:tudien}.

\begin{table}[H]
  \centering
  \caption{Từ điển dữ liệu 10 kênh đặc trưng sinh y tế của tập 548 thuốc}
  \label{tab:tudien}
  \begin{tabularx}{\textwidth}{>{\centering\arraybackslash}p{1.0cm} p{2.6cm} p{2.4cm} X}
    \toprule
    \textbf{STT} & \textbf{Đặc trưng} & \textbf{Kiểu dữ liệu} & \textbf{Ý nghĩa y sinh và phương pháp tính tương đồng} \\
    \midrule
    1 & \texttt{Chem} & SMILES & Cấu trúc phân tử 2D; Morgan Fingerprint, đo Tanimoto \\
    2 & \texttt{Target} & Tập ID Protein & Danh sách protein đích điều trị; đo Jaccard \\
    3 & \texttt{Enzyme} & Tập ID Protein & Enzyme chuyển hóa thuốc trong gan và tế bào; đo Jaccard \\
    4 & \texttt{Transporter} & Tập ID Protein & Protein vận chuyển thuốc qua màng sinh học; đo Jaccard \\
    5 & \texttt{Pathway} & Tập ID Sinh học & Con đường chuyển hóa thuốc (SMPDB/KEGG); đo Jaccard \\
    6 & \texttt{Indication} & Văn bản tự do & Chỉ định lâm sàng của thuốc; TF-IDF + Cosine similarity \\
    7 & \texttt{Side\_effect} & Tập ID UMLS & Các phản ứng phụ thường gặp (SIDER); đo Jaccard \\
    8 & \texttt{Offside\_effect} & Tập ID UMLS & Biến cố bất lợi ngoài nhãn (OFFSIDES); đo Jaccard \\
    9 & \texttt{ATC\_Code} & Mã phân cấp & Mã phân loại giải phẫu, trị liệu của WHO; đo Jaccard \\
    10 & \texttt{MeSH\_Code} & Mã phân cấp & Tiêu đề y khoa chuẩn hóa của NLM; đo Jaccard \\
    \bottomrule
  \end{tabularx}
\end{table}

\subsection{Thống kê mô tả và phân bố bậc tương tác}
\begin{itemize}
  \item Mạng lưới DDI thể hiện rõ tính chất của mạng thế giới nhỏ (scale-free network): Bậc tương tác trung bình của một thuốc là $177{,}31 \pm 129{,}64$, với trung vị là $154{,}0$ tương tác.
  \item Tồn tại một số ``thuốc trung tâm'' (hub drugs) có số lượng tương tác lên tới hơn 400 thuốc khác (như các thuốc an thần, chống co giật), trong khi có những thuốc chỉ tương tác với 1--2 thuốc khác.
\end{itemize}

\subsection{Tương quan liên kênh và mối quan hệ với nhãn}
Qua phân tích ma trận tương quan Pearson giữa 10 nguồn đặc trưng:
\begin{enumerate}
  \item Tất cả hệ số tương quan liên kênh đều nằm dưới $0{,}29$ (mức cao nhất là giữa Target và Chem đạt $r=0{,}283$). Không có hiện tượng đa cộng tuyến ($r > 0{,}85$), khẳng định 10 nguồn phản ánh các khía cạnh sinh học độc lập và bổ khuyết cho nhau hoàn hảo.
  \item Ở các kênh quan trọng như Target, Chem và Indication, độ tương đồng trung bình của nhóm có tương tác ($y=1$) luôn cao hơn rõ rệt (gấp 1,7 đến 1,9 lần) so với nhóm không có tương tác ($y=0$). Hai thuốc cùng gắn kết lên một thụ thể đích hoặc cùng điều trị một chứng bệnh có nguy cơ tương tác đối kháng hoặc cạnh tranh thụ thể rất cao.
\end{enumerate}

% =====================================================================
%  CHƯƠNG 3
% =====================================================================
\chapter{Thực nghiệm và đánh giá}
\label{chap:thucnghiem}

\section{Thiết lập thí nghiệm và các mô hình sử dụng}

\subsection{Môi trường thực nghiệm và giao thức đánh giá}
Các thực nghiệm được thực hiện trên môi trường Python 3.10/3.12 với các thư viện: Scikit-learn, XGBoost, PyTorch Geometric, RDKit, Pandas, NumPy. Toàn bộ các bước có yếu tố ngẫu nhiên đều cố định seed để bảo đảm khả năng tái lập 100\%. Giao thức đánh giá áp dụng phương pháp kiểm định lặp lại qua 10 random seeds độc lập nhằm đo lường độ vững chắc (mean $\pm$ std).

\subsection{Các mô hình sử dụng trong thực nghiệm đối sánh}
Nhóm triển khai và so sánh 5 thuật toán học máy phân loại cùng 1 mô hình tham chiếu cơ sở:
\begin{enumerate}
  \item \textbf{Mô hình cơ sở tham chiếu (GCN Baseline):} Mạng tích chập đồ thị 2 tầng (GCN) kinh điển sử dụng ma trận kề huấn luyện và đặc trưng nút từ ma trận hợp nhất SNF.
  \item \textbf{Hồi quy Logistic (Logistic Regression):} Bộ phân loại tuyến tính với tham số điều chuẩn $\text{max\_iter} = 1000$.
  \item \textbf{K-Nearest Neighbors (KNN):} Thuật toán dựa trên khoảng cách với $K = 5$ láng giềng gần nhất.
  \item \textbf{Rừng ngẫu nhiên (Random Forest):} Thuật toán học tổ hợp đóng bao (Bagging) với $\text{n\_estimators} = 100$ cây quyết định.
  \item \textbf{XGBoost (eXtreme Gradient Boosting):} Thuật toán tăng cường độ dốc mạnh mẽ với $\text{n\_estimators} = 100$, $\text{learning\_rate} = 0{,}1$.
  \item \textbf{Mô hình đề xuất -- Soft Voting Ensemble:} Tích hợp xác suất dự đoán của 4 mô hình (Random Forest, KNN, XGBoost, Logistic Regression) thông qua cơ chế trung bình cộng xác suất với ngưỡng phân loại $0{,}5$:
  \begin{equation}
    P_{\text{ensemble}}(y=1 \mid \mathbf{x}) = \frac{1}{4} \sum_{m=1}^{4} P_m(y=1 \mid \mathbf{x}).
  \end{equation}
  Ưu điểm vượt trội của Soft Voting so với Stacking Ensemble truyền thống là \textbf{không cần huấn luyện thêm một mô hình meta-learner}, do đó loại trừ hoàn toàn nguy cơ rò rỉ dữ liệu ngoài fold (out-of-fold leakage).
\end{enumerate}

\section{Kết quả thực nghiệm và phân tích ý nghĩa}

\subsection{So sánh hiệu năng các mô hình trên 4 kịch bản (Thí nghiệm 1)}
Bảng~\ref{tab:results_main} trình bày kết quả chi tiết của 5 mô hình học máy trên cả 4 kịch bản phân chia dữ liệu, được đo lường bằng 3 chỉ số AUC, AUPR và F1-score theo đúng công bố thực nghiệm của nghiên cứu.

\begin{table}[H]
  \centering
  \caption{So sánh hiệu năng các mô hình học máy trên 4 kịch bản phân chia dữ liệu}
  \label{tab:results_main}
  \begin{tabular}{llccc}
    \toprule
    \textbf{Kịch bản} & \textbf{Mô hình} & \textbf{ROC-AUC} & \textbf{PR-AUC (AUPR)} & \textbf{F1-score} \\
    \midrule
    \multirow{5}{*}{\textbf{Warm-start}}
      & Random Forest       & 0,9299 & 0,9335 & 0,8501 \\
      & Logistic Regression & \textbf{0,9383} & \textbf{0,9441} & \textbf{0,8591} \\
      & KNN                 & 0,7906 & 0,7456 & 0,7153 \\
      & XGBoost             & 0,9300 & 0,9350 & 0,8486 \\
      & \textbf{Soft Voting Ensemble} & 0,9334 & 0,9373 & 0,8519 \\
    \midrule
    \multirow{5}{*}{\textbf{Strict Cold-start}}
      & Random Forest       & 0,6674 & 0,6512 & 0,5567 \\
      & Logistic Regression & 0,6839 & 0,6671 & 0,4948 \\
      & KNN                 & 0,5457 & 0,5283 & \textbf{0,6351} \\
      & XGBoost             & \textbf{0,6900} & \textbf{0,6884} & 0,5993 \\
      & \textbf{Soft Voting Ensemble} & 0,6828 & 0,6791 & 0,6064 \\
    \midrule
    \multirow{5}{*}{\textbf{Semi Cold-start}}
      & Random Forest       & 0,8310 & 0,8183 & 0,7499 \\
      & Logistic Regression & 0,8350 & 0,8226 & 0,7296 \\
      & KNN                 & 0,7833 & 0,7425 & 0,7135 \\
      & XGBoost             & 0,8260 & 0,8270 & 0,7337 \\
      & \textbf{Soft Voting Ensemble} & \textbf{0,8428} & \textbf{0,8313} & \textbf{0,7521} \\
    \midrule
    \multirow{5}{*}{\textbf{Mixed Scenario}}
      & Random Forest       & 0,8175 & 0,8059 & 0,7271 \\
      & Logistic Regression & 0,8228 & 0,8115 & 0,7091 \\
      & KNN                 & 0,7627 & 0,7209 & 0,7034 \\
      & XGBoost             & 0,8134 & 0,8089 & 0,7205 \\
      & \textbf{Soft Voting Ensemble} & \textbf{0,8302} & \textbf{0,8202} & \textbf{0,7375} \\
    \bottomrule
  \end{tabular}
\end{table}

\textit{Nhận xét và phân tích ý nghĩa khoa học từ Bảng~\ref{tab:results_main}:}
\begin{enumerate}
  \item \textbf{Ở kịch bản Warm-start:} Logistic Regression đạt hiệu năng cao nhất (AUC 0,9383, AUPR 0,9441, F1 0,8591). Khi cấu trúc đồ thị quen thuộc có mặt đầy đủ ở cả hai đầu mút, ranh giới quyết định được phân tách rất tốt và mô hình tuyến tính đơn giản là đủ để đạt kết quả xuất sắc mà không cần độ phức tạp phi lý.
  \item \textbf{Ở kịch bản khắt khe Strict Cold-start (New $\times$ New):} XGBoost thể hiện ưu thế vượt trội về phân loại phi tuyến tính (AUC 0,6900, AUPR 0,6884), trong khi KNN đạt F1-score cao nhất (0,6351). Hiệu năng của tất cả các mô hình đều sụt giảm rõ rệt so với Warm-start, chứng minh sự thật khách quan rằng: \textit{Random split trước đây đã mang lại kết quả quá lạc quan do rò rỉ dữ liệu.}
  \item \textbf{Sự ưu việt và tính ổn định của Soft Voting Ensemble:} Ở hai kịch bản thực tế nhất là Semi Cold-start và Mixed, mô hình Ensemble đề xuất giành chiến thắng toàn diện trên cả 3 chỉ số (Semi: AUC 0,8428, AUPR 0,8313, F1 0,7521; Mixed: AUC 0,8302, AUPR 0,8202, F1 0,7375). Việc dung hòa các quyết định dự đoán giúp mô hình thích ứng bền bỉ nhất trước sự biến động dữ liệu lâm sàng.
\end{enumerate}

\subsection{Đánh giá độ vững chắc và đối sánh với GCN Baseline qua 10 Seeds (Thí nghiệm 2)}
Để chứng minh tính vững chắc khoa học và vạch trần lỗ hổng của các mô hình học sâu đồ thị (GNNs) khi gặp bài toán khởi động lạnh, nhóm thực hiện thực nghiệm đối sánh lặp lại qua 10 random seeds độc lập giữa mô hình Soft Voting Ensemble và mô hình GCN Baseline (Bảng~\ref{tab:robustness}).

\begin{table}[H]
  \centering
  \caption{Đánh giá độ vững chắc (Chỉ số AUPR qua 10 random seeds: Mean $\pm$ Std)}
  \label{tab:robustness}
  \begin{tabularx}{\textwidth}{l >{\centering\arraybackslash}X >{\centering\arraybackslash}X >{\centering\arraybackslash}X >{\centering\arraybackslash}X}
    \toprule
    \textbf{Mô hình đối sánh} & \textbf{Case 1 (Warm)} & \textbf{Case 2 (Strict)} & \textbf{Case 3 (Semi)} & \textbf{Case 4 (Mixed)} \\
    \midrule
    GCN Baseline     & $0{,}9356 \pm 0{,}0018$ & $0{,}4695 \pm 0{,}0341$ & $0{,}7733 \pm 0{,}0136$ & $0{,}7560 \pm 0{,}0199$ \\
    \textbf{Ensemble (Đề xuất)} & $\mathbf{0{,}9363 \pm 0{,}0018}$ & $\mathbf{0{,}6720 \pm 0{,}0247}$ & $\mathbf{0{,}8195 \pm 0{,}0116}$ & $\mathbf{0{,}8079 \pm 0{,}0118}$ \\
    \bottomrule
  \end{tabularx}
\end{table}

\textit{Nhận xét đột phá về mặt khoa học:}
\begin{itemize}
  \item Ở Case 1 (Warm-start), GCN hoạt động rất ấn tượng ($0{,}9356$) nhờ khai thác triệt để các liên kết cạnh có sẵn.
  \item Tuy nhiên, khi chuyển sang Case 2 (Strict Cold-start), hiệu năng của GCN \textbf{sụt giảm nghiêm trọng xuống chỉ còn $0{,}4695$} (thấp hơn cả mức phân loại ngẫu nhiên $0{,}50$). Nguyên nhân là do các thuốc mới bị cô lập hoàn toàn, không có cạnh nào kết nối vào đồ thị huấn luyện, khiến cơ chế truyền tin (message passing) của mạng nơ-ron đồ thị bị vô hiệu hóa hoàn toàn.
  \item Ngược lại, mô hình Soft Voting Ensemble của nhóm vẫn duy trì AUPR đạt $0{,}6720$, vượt trội hơn GCN tới hơn 20 điểm phần trăm, chứng minh giá trị to lớn của việc kết hợp các góc nhìn đặc trưng sinh hóa độc lập.
\end{itemize}

\subsection{Phân tích lỗi định tính (Qualitative Error Analysis)}
Nhóm tiến hành phân tích sâu các mẫu dương tính giả (False Positives -- các cặp được mô hình dự đoán có tương tác với xác suất rất cao $p > 0{,}8$ nhưng chưa có nhãn trong DrugBank ở Case 2). Trong dược lý học, do các cơ sở dữ liệu y tế luôn tồn tại độ trễ cập nhật (Positive-Unlabeled data), các cặp thuốc này có độ tương đồng cực cao về cấu trúc phân tử và protein đích tác dụng. Nhiều cặp trong số này hoàn toàn có thể là những \textit{tương tác thuốc tiềm ẩn có thật trong thực tế} mà y văn chưa kịp công bố, khẳng định tiềm năng ứng dụng của mô hình như một công cụ hỗ trợ sàng lọc tiền lâm sàng hữu ích.

\subsection{Thực nghiệm bóc tách thành phần (Ablation Study -- Thí nghiệm 4)}
Để kiểm tra mức độ đóng góp độc lập của từng khối thành phần trong pipeline xử lý, nhóm thực hiện thử nghiệm loại bỏ lần lượt từng thành phần (Bảng~\ref{tab:ablation}).

\begin{table}[H]
  \centering
  \caption{Kết quả bóc tách thành phần (Chỉ số AUPR khi loại bỏ các module)}
  \label{tab:ablation}
  \begin{tabularx}{\textwidth}{l >{\centering\arraybackslash}X >{\centering\arraybackslash}X >{\centering\arraybackslash}X >{\centering\arraybackslash}X}
    \toprule
    \textbf{Cấu hình thực nghiệm} & \textbf{Case 1 (Warm)} & \textbf{Case 2 (Strict)} & \textbf{Case 3 (Semi)} & \textbf{Case 4 (Mixed)} \\
    \midrule
    \textbf{Full Pipeline (Đầy đủ)} & \textbf{0,9373} & 0,6791 & 0,8313 & 0,8202 \\
    Without Chem (Bỏ cấu trúc SMILES) & 0,9371 & 0,6700 & 0,8295 & 0,8181 \\
    Without Targets (Bỏ protein đích) & 0,9363 & 0,6904 & 0,8245 & 0,8144 \\
    Without SNF (Bỏ hợp nhất mạng)    & 0,9363 & \textbf{0,7010} & \textbf{0,8371} & \textbf{0,8264} \\
    \bottomrule
  \end{tabularx}
\end{table}

\textit{Phát hiện thú vị và trung thực về mặt khoa học:}
\begin{itemize}
  \item Việc loại bỏ đặc trưng hóa học (\texttt{Without Chem}) làm suy giảm AUPR trên toàn bộ cả 4 kịch bản, khẳng định vai trò nền tảng không thể thay thế của cấu trúc phân tử 2D.
  \item Đáng chú ý, việc loại bỏ hợp nhất mạng SNF (\texttt{Without SNF}) lại làm tăng nhẹ chỉ số AUPR ở các kịch bản Cold-start (Case 2, 3, 4). Phát hiện trung thực này mang lại đóng góp học thuật sâu sắc: \textit{Quá trình hợp nhất mạng tương đồng (SNF) dù làm giàu thông tin trong môi trường transductive nhưng có thể vô tình đưa thêm nhiễu (noise) vào không gian biểu diễn khi gặp các thực thể thuốc hoàn toàn mới.}
\end{itemize}

% =====================================================================
%  KẾT LUẬN
% =====================================================================
\unnumberedchapter{KẾT LUẬN VÀ HƯỚNG PHÁT TRIỂN}

\section*{1. Các kết quả chính đạt được của đề tài}
Sau quá trình nghiên cứu và thực nghiệm nghiêm túc, đề tài đã giải quyết trọn vẹn các mục tiêu đặt ra:
\begin{enumerate}
  \item \textbf{Khắc phục triệt để hiện tượng rò rỉ dữ liệu:} Đề xuất thành công khung đánh giá 4 kịch bản phân tách nghiêm ngặt thực thể thuốc, chứng minh rõ ràng mức độ suy giảm hiệu năng thực tế khi chuyển từ Warm-start sang Cold-start.
  \item \textbf{Xây dựng không gian đặc trưng đa nguồn chất lượng cao:} Tích hợp 10 kênh đặc trưng sinh y tế từ 548 thuốc DrugBank 5.1.13, áp dụng lọc Entropy và hợp nhất SNF bài bản.
  \item \textbf{Mô hình Soft Voting Ensemble hiệu quả và ổn định:} Duy trì khả năng tổng quát hóa vượt trội ở các kịch bản thực tế (Semi Cold đạt AUPR 0,8313; Mixed đạt AUPR 0,8202), khắc phục hoàn toàn sự sụp đổ nghiêm trọng của mô hình GCN Baseline ở kịch bản Strict Cold-start.
  \item \textbf{Đóng góp học thuật minh bạch:} Phân tích trung thực hiện tượng nhiễu của SNF trong điều kiện khởi động lạnh thông qua Ablation Study.
\end{enumerate}

\section*{2. Những hạn chế còn tồn tại}
\begin{itemize}
  \item Nghiên cứu hiện mới dừng lại ở bài toán phân loại nhị phân (có/không có DDI), chưa phân loại cụ thể mức độ nghiêm trọng hay loại biến cố bất lợi cụ thể (multi-typed DDI).
  \item Quy mô dữ liệu hiện được kiểm nghiệm trên tập chuẩn 548 thuốc; việc mở rộng lên toàn bộ hơn 2.000 thuốc trong DrugBank cần thêm tài nguyên phần cứng lớn.
\end{itemize}

\section*{3. Hướng phát triển trong tương lai}
\begin{enumerate}
  \item Phát triển các kiến trúc mạng nơ-ron đồ thị cải tiến có khả năng suy luận quy nạp (Inductive GNNs) kết hợp cơ chế chú ý (Graph Attention) để tận dụng cấu trúc mạng mà không bị sụp đổ ở cold-start.
  \item Tích hợp thêm các mô hình ngôn ngữ lớn y sinh (BioBERT / PubMedBERT) để trích xuất tự động thông tin tương tác từ các bài báo y khoa cập nhật thời gian thực.
\end{enumerate}

% =====================================================================
%  TÀI LIỆU THAM KHẢO
% =====================================================================
\begin{thebibliography}{99}
\setlength{\itemsep}{2pt}

\bibitem{mitchell1997}
T.~M.~Mitchell (1997), \textit{Machine Learning}, McGraw-Hill, New York.

\bibitem{pirmohamed2004adverse}
M.~Pirmohamed \textit{et al.} (2004), ``Adverse drug reactions as cause of admission to hospital: prospective analysis of 18 820 patients'', \textit{BMJ}, vol.~329, no.~7456, pp.~15--19.

\bibitem{vilar2014}
S.~Vilar \textit{et al.} (2014), ``Similarity-based modeling for predicting drug-drug interactions'', \textit{Nature Protocols}, vol.~9, no.~9, pp.~2147--2163.

\bibitem{rohani2019}
N.~Rohani, H.~Eslahchi (2019), ``Drug-Drug Interaction Prediction by Integrating Multiple Biological Data Sources'', \textit{IEEE/ACM Transactions on Computational Biology and Bioinformatics}, vol.~17, no.~6, pp.~1896--1908.

\bibitem{zitnik2018}
M.~Zitnik, M.~Agrawal, J.~Leskovec (2018), ``Modeling polypharmacy side effects with graph convolutional networks'', \textit{Bioinformatics}, vol.~34, no.~13, pp.~i457--i466.

\bibitem{wishart2018}
D.~S.~Wishart \textit{et al.} (2018), ``DrugBank 5.0: a major update to the DrugBank database for 2018'', \textit{Nucleic Acids Research}, vol.~46, no.~D1, pp.~D1074--D1082.

\bibitem{kuhn2016}
M.~Kuhn \textit{et al.} (2016), ``The SIDER database of drugs and side effects'', \textit{Nucleic Acids Research}, vol.~44, no.~D1, pp.~D1075--D1079.

\bibitem{tatonetti2012}
N.~P.~Tatonetti \textit{et al.} (2012), ``Data-driven prediction of drug effects and interactions'', \textit{Science Translational Medicine}, vol.~4, no.~125, pp.~125ra31.

\bibitem{wang2014}
B.~Wang \textit{et al.} (2014), ``Similarity network fusion for aggregating data types on a genomic scale'', \textit{Nature Methods}, vol.~11, no.~3, pp.~333--337.

\bibitem{breiman2001}
L.~Breiman (2001), ``Random forests'', \textit{Machine Learning}, vol.~45, no.~1, pp.~5--32.

\bibitem{chen2016}
T.~Chen, C.~Guestrin (2016), ``XGBoost: A scalable tree boosting system'', in \textit{Proceedings of the 22nd ACM SIGKDD International Conference on Knowledge Discovery and Data Mining}, pp.~785--794.

\bibitem{rogers2010}
D.~Rogers, M.~Hahn (2010), ``Extended-connectivity fingerprints'', \textit{Journal of Chemical Information and Modeling}, vol.~50, no.~5, pp.~742--754.

\bibitem{langground2026}
A.~Ajit (2026), ``Language-Grounded Drug Representations for Cold-Start Drug-Drug Interaction Prediction'', \textit{Preprint}, Research Square, doi: 10.21203/rs.3.rs-10334947/v1.

\bibitem{fgddi2026}
F.~Zhou, S.~Uddin (2026), ``FG-DDI: Functional group-aware graph neural networks for drug-drug interaction prediction'', \textit{Journal of Biomedical Informatics}, vol.~174, p.~104981.

\bibitem{mrlfddi2026}
J.~Zhong \textit{et al.} (2026), ``MRLF-DDI: A Multi-View Representation Learning Framework for Drug-Drug Interaction Event Prediction'', \textit{IEEE Journal of Biomedical and Health Informatics}, vol.~30, no.~3, pp.~2217--2227.

\end{thebibliography}

% =====================================================================
%  PHỤ LỤC
% =====================================================================
\appendix
\renewcommand{\bmchapname}{Phụ lục}
\addtocontents{toc}{\protect\renewcommand{\protect\tocchapname}{\appendixname}}

\chapter{Hướng dẫn cài đặt và chạy mã nguồn thực nghiệm}
\label{pl:huongdan}

\section{Yêu cầu môi trường thực thi}
\begin{itemize}
  \item Python phiên bản $\ge 3.9$ (khuyến nghị Python 3.10 hoặc 3.11).
  \item Các gói thư viện: \texttt{pandas}, \texttt{numpy}, \texttt{scikit-learn}, \texttt{xgboost}, \texttt{rdkit}, \texttt{matplotlib}, \texttt{seaborn}, \texttt{torch}, \texttt{torch-geometric}.
\end{itemize}

\section{Quy trình tái lập kết quả}
\begin{enumerate}
  \item \textbf{Cài đặt môi trường:}
\begin{lstlisting}[language=bash,numbers=none]
pip install pandas numpy scikit-learn xgboost rdkit matplotlib seaborn torch
\end{lstlisting}
  \item \textbf{Thực thi phân tích dữ liệu và trích xuất ma trận:}
\begin{lstlisting}[language=bash,numbers=none]
python generate_eda_and_figures.py
\end{lstlisting}
  \item \textbf{Chạy thực nghiệm mô hình và đánh giá 4 kịch bản:}
  Khởi chạy notebook \texttt{test\_case\_1\_random\_split.ipynb} hoặc file thực nghiệm tổng hợp trong thư mục \texttt{code\_4\_case/}. Toàn bộ bảng kết quả Table 2, Table 3 và Table 4 sẽ được tự động tính toán và kết xuất chính xác.
\end{enumerate}

\chapter{Nhật ký công việc nhóm theo tuần}
\label{pl:nhatky}

\begin{table}[H]
  \centering
  \caption{Nhật ký công việc nhóm từ Tuần 1 đến Tuần 6}
  \label{tab:nhatky}
  \begin{tabularx}{\textwidth}{>{\centering\arraybackslash}p{1.2cm} >{\centering\arraybackslash}p{2.6cm} p{3.2cm} X}
    \toprule
    \textbf{Tuần} & \textbf{Thời gian} & \textbf{Thành viên} & \textbf{Nội dung công việc và kết quả hoàn thành} \\
    \midrule
    \multirow{2}{*}{1} & \multirow{2}{*}{01/10 -- 07/10/2026} & \svAten & Tiếp nhận đề tài; khảo sát tài liệu tổng quan DDI; phân tích hạn chế của chia ngẫu nhiên. \\ \cline{3-4}
                       &                                       & \svBten & Cài đặt môi trường; khảo sát các nghiên cứu liên quan (Rohani 2019, Vilar 2014, Zitnik 2018). \\ \midrule
    \multirow{2}{*}{2} & \multirow{2}{*}{08/10 -- 14/10/2026} & \svAten & Thu thập tập dữ liệu 548 thuốc; trích xuất 10 kênh đặc trưng sinh học; xử lý dữ liệu thiếu. \\ \cline{3-4}
                       &                                       & \svBten & Xây dựng pipeline tính ma trận tương đồng Tanimoto, Cosine và Jaccard cho 10 kênh. \\ \midrule
    \multirow{2}{*}{3} & \multirow{2}{*}{15/10 -- 21/10/2026} & \svAten & Xây dựng module lọc Entropy và thuật toán hợp nhất mạng tương đồng (SNF, K=20, t=20). \\ \cline{3-4}
                       &                                       & \svBten & Thực hiện phân tích EDA; đo lường tương quan liên kênh và phân bố bậc thuốc trung tâm. \\ \midrule
    \multirow{2}{*}{4} & \multirow{2}{*}{22/10 -- 28/10/2026} & \svAten & Thiết lập 4 kịch bản chia tập chống rò rỉ (Warm, Strict Cold, Semi Cold, Mixed). \\ \cline{3-4}
                       &                                       & \svBten & Lấy mẫu âm cân bằng 1:1; hoàn thiện báo cáo tiểu luận Bài kiểm tra 1 (Giai đoạn 1). \\ \midrule
    \multirow{2}{*}{5} & \multirow{2}{*}{29/10 -- 04/11/2026} & \svAten & Xây dựng mô hình Soft Voting Ensemble tích hợp RF, KNN, XGBoost và Logistic Regression. \\ \cline{3-4}
                       &                                       & \svBten & Cài đặt mô hình GCN Baseline trên PyTorch Geometric để làm mốc đối sánh. \\ \midrule
    \multirow{2}{*}{6} & \multirow{2}{*}{05/11 -- 11/11/2026} & \svAten & Tổng hợp kết quả Table 2, thực hiện Ablation Study (Table 4) và viết hoàn chỉnh báo cáo. \\ \cline{3-4}
                       &                                       & \svBten & Chạy 10-seed đánh giá độ vững chắc (Table 3), phân tích lỗi False Positives và nộp báo cáo. \\
    \bottomrule
  \end{tabularx}
\end{table}

\end{document}
